\clearpage
\SourceLeaf{479}
\pdfbookmark[1]{Nachlass III. Specielle Störungen der Pallas durch Jupiter. Erste Rechnung.}{pallas-III}
\fancyhead[L]{\small NACHLASS. STÖRUNGEN DER PALLAS}
\begin{center}
{\Large NACHLASS.\par}
\bigskip
{\large [III. SPECIELLE STÖRUNGEN DER PALLAS DURCH JUPITER. ERSTE RECHNUNG.]\par}
\medskip
[OCTOBER—DECEMBER 1810.]
\end{center}
\par\Needspace{9\baselineskip}\bigskip
\pdfbookmark[2]{Erste Rechnung, [1]}{pallas-III-fragment1}
\begin{center}[1.]\end{center}
\begingroup
\makeatletter\tagsleft@true\let\veqno\@@leqno\makeatother
\setlength{\lineskip}{4pt}\setlength{\lineskiplimit}{4pt}
\noindent\hspace*{\parindent}[Die Rechnung bezieht sich auf die Störungen der Elemente und erstreckt sich über den 3000 Tage umfassenden Zeitraum von 1803 Juni 30 bis 1811 September 16, welcher in 60 Intervalle zu je 50 Tagen eingetheilt wurde, so dass die Rechnungen im Ganzen für 61 Zeitmomente auszuführen waren. Sie stützt sich im Wesentlichen auf Formeln, die mit unerheblichen Modificationen in der »Exposition d’une nouvelle méthode« etc. abgeleitet sind; vgl. Art. 20 dieser Abhandlung (Seite 464 dieses Bandes).
\par
Unter Annahme der Ebene der Pallasbahn als Grundebene wurden die folgenden Bezeichnungen benutzt:
\begin{center}
\begin{tabular}{@{}r@{\ }p{.86\linewidth}@{}}
\(r={}\)&Radius Vector der Pallas\\
\(v={}\)&heliocentrische Länge der Pallas in ihrer Bahn\\
\(E={}\)&excentrische Anomalie der Pallas\\
\(r'={}\)&Radius Vector des Jupiter\\
\(w'={}\)&heliocentrische Länge des Jupiter in seiner Bahn\\
\(v'={}\)&Länge des Jupiter, projicirt auf die Ebene der Pallasbahn\\
\(\beta'={}\)&Breite des Jupiter über der Ebene der Pallasbahn\\
\(\rho={}\)&gegenseitige Entfernung von Pallas und Jupiter\\
\(\mu={}\)&Jupitersmasse in Einheiten der Sonnenmasse\\
\(k={}\)&Constante der Theoria motus
\end{tabular}
\end{center}
\[
\begin{gathered}
\left.
\begin{aligned}
\begin{gathered}a,\ e,\ \varphi,\ \tilde\omega,\ p=a(1-ee),\\ n,\ \varepsilon,\ i,\ \AscendingNode\end{gathered}&=\text{Elemente der Pallasbahn}\\
a',\ e'\ \text{etc.}&=\text{Elemente der Jupitersbahn}
\end{aligned}
\right\}\quad\\
\text{in ihrer gewöhnlichen Bedeutung.}
\end{gathered}
\]
\par
Wenn man zur Abkürzung setzt:
\[
\begin{aligned}
r'\cos\beta'\cos(v-v')&=\xi&\qquad\frac{1}{\rho^{3}}-\frac{1}{r'^{3}}&=\Delta\\
r'\cos\beta'\sin(v-v')&=\eta&&\\
r'\sin\beta'&=\zeta&\frac{\mu k\mathrm{d}t}{\BareRoot{p}}&=(1),
\end{aligned}
\tag*{1)}
\]
\SourceLeaf{480}womit auch \(rr+r'r'-2r\xi=\rho\rho\) wird, so berechnen sich die Grössen \(T,V,W\), welche sich hier von den gleichbezeichneten Grössen in der Exposition Art. 15 (Seite 458) durch die Factoren \(-\dfrac{a}{\cos\varphi}\) resp. \(-\dfrac{ar}{\cos\varphi}\) resp. \(+\dfrac{ar}{\cos\varphi}\) unterscheiden, nach den Formeln:
\[
T=\frac{(1)r}{\rho^{3}}-(1)\Delta\xi,\qquad V=(1)r\Delta\eta,\qquad W=(1)r\Delta\zeta,
\tag*{2)}
\]
und hieraus finden sich die Beträge der momentanen Störungen während des Zeitraums \(\mathrm{d}t\):
\[
\begin{aligned}
\mathrm{d}i&=\cos(v-\AscendingNode).W\\
\mathrm{d}\AscendingNode&=(2)\sin(v-\AscendingNode).W\\
\mathrm{d}n&=(4)\sin(v-\tilde\omega).T+\frac{(5)}{rr}V\\
\mathrm{d}\varphi&=(6)\sin(v-\tilde\omega).T+\frac{(6)(\cos(v-\tilde\omega)+\cos E)}{r}V\\
\mathrm{d}\tilde\omega&=A+(7)\cos(v-\tilde\omega).T\\
&\qquad+(8)\sin(\dot v-\tilde\omega).\left(1+\frac{p}{r}\right)V\\
\mathrm{d}\varepsilon&=A+(9)rT+(10)\cos(v-\tilde\omega).T\\
&\qquad+(11)\sin(v-\tilde\omega).\left(1+\frac{p}{r}\right)V,
\end{aligned}
\tag*{3)}
\]
wo der Kürze halber bezeichnet ist
\[
A=(3)\sin(v-\AscendingNode).W
\]
\[
\begin{aligned}
(2)&=\frac{1}{\sin i}&\qquad(7)&=\frac{p}{e}\\[3pt]
(3)&=\operatorname{tang}\tfrac{1}{2}i&(8)&=-\frac{1}{e}\\[3pt]
(4)&=3aen&(9)&=2\cos\varphi\\[3pt]
(5)&=3apn&(10)&=a\cos\varphi^{2}\operatorname{tang}\tfrac{1}{2}\varphi\\[3pt]
(6)&=-a\cos\varphi&(11)&=-\operatorname{tang}\tfrac{1}{2}\varphi.]
\end{aligned}
\]
\par\endgroup
